\title{Controlling Text-to-Image Diffusion by Orthogonal Finetuning}

\begin{abstract}
Recent advances in large text-to-image diffusion models have demonstrated substantial success in producing highly realistic images conditioned on textual descriptions. However, effectively steering these sophisticated models to fulfill various downstream objectives remains a significant unresolved issue. In response, we propose a systematic finetuning approach termed Orthogonal Finetuning (OFT), designed for adapting text-to-image diffusion models to specialized tasks. Unlike prior techniques, OFT is theoretically guaranteed to maintain hyperspherical energy, which encodes the pairwise interactions among neurons when projected onto the unit hypersphere. Our analysis suggests that conserving this quantity is essential for retaining the semantic fidelity of the generated imagery. To further enhance the robustness of the finetuning process, we introduce Constrained Orthogonal Finetuning (COFT), which augments OFT by enforcing a radius constraint within the hyperspherical space. We apply these methodologies to two central finetuning scenarios: subject-driven generation—where the challenge lies in synthesizing diverse depictions of a particular subject from a limited image set alongside a text prompt—and controllable generation—where the model is conditioned on supplementary control inputs. Our empirical studies demonstrate that the OFT method achieves superior generation quality and faster convergence rates compared to existing approaches.
\end{abstract}

\section{Introduction}

State-of-the-art text-to-image diffusion models~\cite{saharia2022photorealistic,ramesh2022hierarchical,rombach2022high} have recently set new standards in synthesizing visually realistic images based on text prompts. While their performance is noteworthy, there remain notable challenges: text-based conditioning can often lack precision and granularity, thereby restricting exact control over the generated outputs. To confront these issues, this paper focuses on two critical variants of text-to-image generation tasks:

\begin{itemize}[leftmargin=*,nosep]

    \item \textbf{Subject-driven generation}~\cite{ruiz2023dreambooth}: With access to only a handful of images of a target subject, the objective is to produce new images portraying that subject in novel contexts under textual guidance.
    \item \textbf{Controllable generation}~\cite{zhang2023adding,mou2023t2i}: Provided with extra control inputs (such as canny edge maps or segmentation masks), the aim is to generate images that conform to both the external control signal and the descriptive text prompt.
\end{itemize}

At their core, both of these tasks revolve around the challenge of adapting text-to-image diffusion models through finetuning, while ensuring that the high-quality generative capabilities established during pretraining are not compromised. The criteria for a robust finetuning approach can be outlined as follows: (1) \emph{training efficiency}, which involves reducing the number of trainable parameters and minimizing training epochs; and (2) \emph{generalizability preservation}, ensuring the retention of both fidelity and diversity in generated outputs. Typical finetuning approaches proceed by either adjusting the model’s weights using a low learning rate (\emph{e.g.}, \cite{ruiz2023dreambooth}), or by introducing minor architectural modifications with re-parameterized weights (\emph{e.g.}, \cite{hulora2022,zhang2023adding}). However, while these methods are straightforward, neither inherently ensures that the pretrained generative strengths are maintained. Moreover, there remains a limited theoretical framework to guide the construction of effective finetuning methodologies or the selection of suitable hyperparameters such as training duration. A fundamental obstacle is the absence of a reliable metric for evaluating how well pretrained generative capabilities are conserved post-finetuning. Current techniques tend to operate under the assumption that a reduced Euclidean distance between the finetuned and original model correlates with greater preservation of generative ability. Consequently, finetuning is often performed with exceedingly low learning rates. Nonetheless, this assumption is not universally valid. We contend that Euclidean proximity to the pretrained model fails to fully capture the extent of semantic retention. Instead, more structurally-informed metrics for assessing differences between finetuned and pretrained models may provide substantive benefits, improving both the stability of finetuning and the safeguarding of pretrained generative performance.

Drawing upon empirical findings that hyperspherical similarity effectively captures semantic relationships~\cite{liu2017deep,liu2018decoupled,chen2020angular}, our approach leverages hyperspherical energy~\cite{liu2018learning} to quantify the pairwise dependencies among neurons. Specifically, hyperspherical energy is computed as the aggregate of hyperspherical similarities (such as cosine similarity) among all neuron pairs within a layer, thereby measuring the extent to which neuron activations are uniformly distributed on the unit hypersphere~\cite{liu2021learning}. Our conjecture is that a well-adapted finetuned model should exhibit only marginal change in hyperspherical energy relative to its pretrained counterpart. While one straightforward strategy is to regularize training so that hyperspherical energy is conserved during finetuning, this approach does not inherently guarantee minimization of the energy gap. To address this, we exploit a key invariance property: hyperspherical energy, and thus all pairwise hyperspherical similarities, remain unchanged under any global orthogonal transformation applied to all neurons. Building on this, we introduce Orthogonal Finetuning~(OFT), a method designed to adapt large text-to-image diffusion models for downstream tasks while strictly preserving their hyperspherical energy. The principle behind OFT is to optimize a single orthogonal transformation, shared across each layer, ensuring the pairwise angles between neurons are unaltered. This procedure can also be conceptualized as recalibrating the canonical basis for the set of neurons within a layer. Furthermore, motivated by the observation that reduced Euclidean distance between pretrained and finetuned model weights tends to safeguard the original performance, we put forth a refinement—Constrained Orthogonal Finetuning~(COFT)—which restricts the finetuned model to remain within a fixed-radius hypersphere centered at the pretrained neuron weights.

The rationale behind employing orthogonal transformations in neuron finetuning is partly inspired by concepts from the 2D Fourier transform, wherein an image is represented by its magnitude and phase spectra. Notably, the phase spectrum encodes angular relationships between the input and the basis, encapsulating most of the semantic information. For instance, reconstructing an image using its phase spectrum combined with an arbitrary magnitude spectrum can still yield the semantic content of the original image. This observation implies that manipulating the orientation of neurons is fundamental for semantic alteration in generated images, a central challenge in both subject-driven and controllable generation. Nevertheless, unconstrained modifications to neuron orientations can undermine the capabilities inherited from pretraining. To address this, we advocate for maintaining the pairwise angular relationships among neurons, grounded in the assumption that such angles play a pivotal role in encoding network knowledge. Consequently, we propose to apply a shared orthogonal transformation within each layer, thereby keeping the hyperspherical energy invariant.

Additionally, our method is influenced by the work on orthogonal over-parameterized training~\cite{liu2021orthogonal}, where classification networks are trained from a random initialization by orthogonally transforming the initial weights. This is based on the property that a randomly initialized network tends to have low hyperspherical energy on average, and the objective in~\cite{liu2021orthogonal} is to maintain low energy during training (as reduced energy has been shown to enhance generalization performance in classification~\cite{liu2018learning,lin2020regularizing}). Their findings indicate that orthogonal transformations offer enough capacity for training robust networks in classification tasks. By contrast, our focus is on the finetuning of text-to-image diffusion models to enable greater control and improved generative results. There are two key distinctions between OFT and the method in~\cite{liu2021orthogonal}. Firstly, while~\cite{liu2021orthogonal} aims to minimize hyperspherical energy, OFT seeks to conserve the hyperspherical energy of the pretrained model, preserving the model’s original structure during finetuning; in fact, reducing hyperspherical energy during diffusion model finetuning can damage established semantic structures. Secondly, OFT is formulated to minimize changes relative to the pretrained network, introducing an explicit constraint absent in~\cite{liu2021orthogonal}. Achieving an optimal balance between adaptability and preservation of pretrained knowledge is critical in finetuning, and we contend that OFT successfully navigates this trade-off. Our main contributions are summarized as follows:

\begin{itemize}[leftmargin=*,nosep]

    \item We introduce a new finetuning technique, Orthogonal Finetuning, which enhances the controllability of text-to-image diffusion models. To promote greater training robustness, we develop a constrained variant that restricts angular divergence from the original pretrained weights.
    \item In contrast with traditional finetuning methods, OFT preserves the hyperspherical energy throughout model adaptation—a property that we theoretically ensure and empirically observe to be crucial for maintaining the pretrained model’s generative semantics.
    \item Our method, OFT, is evaluated on subject-driven and controllable image generation tasks. Extensive experiments reveal that it significantly outperforms previous approaches in terms of image quality, training efficiency, and stability during finetuning. Notably, OFT is more sample-efficient, requiring fewer images and epochs to reach strong performance.
    \item For evaluating controllable generation, we propose a novel control consistency metric. The approach assesses controllability by estimating the control signal from the synthesized image and comparing it to the original input control, further substantiating the high controllability achievable by OFT.
\end{itemize}

\section{Related Work}

\textbf{Text-to-image diffusion models}. The field of text-to-image generation has rapidly advanced~\cite{saharia2022photorealistic,ramesh2022hierarchical,rombach2022high,gu2022vector,nichol2022glide}, propelled by progress in diffusion generative frameworks~\cite{ho2020denoising,song2021score,song2021denoising,dhariwal2021diffusion} and joint vision-language learning~\cite{su2020vlbert,lu2019vilbert,radford2021learning,wang2021simvlm,li2022blip,li2023blip,alayrac2022flamingo,schuhmann2022laion}. Models such as GLIDE~\cite{nichol2022glide} and Imagen~\cite{saharia2022photorealistic} perform diffusion in pixel space, with GLIDE employing joint training of the text encoder and diffusion prior on text-image pairs, and Imagen leveraging a fixed pretrained text encoder. Meanwhile, Stable Diffusion~\cite{rombach2022high} and DALL-E2~\cite{ramesh2022hierarchical} operate in a learned latent space; Stable Diffusion utilizes VQ-GAN~\cite{esser2021taming} for its latent visual codebook, while DALL-E2 uses CLIP~\cite{radford2021learning} to build a shared latent embedding space for text and images. Alongside diffusion models, alternative approaches—such as generative adversarial networks~\cite{reed2016generative,zhang2017stackgan,xu2018attngan,li2019controllable} and autoregressive strategies~\cite{ramesh2021zero,ding2021cogview,wu2022nuwa,yuscaling}—have also been explored for text-conditional image generation. OFT is a generic, model-agnostic finetuning strategy that is compatible with any text-to-image diffusion model.

\textbf{Subject-driven generation}. To avoid altering the primary subject in generated images, previous methods~\cite{avrahami2022blended,nichol2022glide} leverage user-supplied masks as conditional inputs. Inversion-based approaches~\cite{choi2021ilvr,dhariwal2021diffusion,ramesh2022hierarchical,gal2022image} are employed to manipulate image backgrounds or context without changing the subject. The method in~\cite{hertz2022prompt} enables both local and global editing absent any explicit masking. These approaches, however, often struggle to retain the nuanced identity-specific features of the subject. Pivotal Tuning~\cite{roich2022pivotal} finetunes a generator in the vicinity of a starting latent inversion and adds regularization to help maintain identity. Likewise, the method in~\cite{nitzan2022mystyle} creates a personalized generative prior using a set of facial images belonging to a single individual. The work by~\cite{casanova2021instance} generates multiple variations for a given instance but might compromise instance-specific information. DreamBooth~\cite{ruiz2023dreambooth} improves the diffusion model for a novel subject by finetuning it with a unique token and a handful of subject images, employing both reconstruction and class-specific prior preservation losses. In our approach, OFT adopts the overall DreamBooth setup; rather than simple finetuning with a low learning rate, we employ orthogonal transformations during the finetuning process.

\textbf{Controllable generation}. Image-to-image translation is a specific case within the broader domain of controllable generation. Prior approaches predominantly leverage conditional generative adversarial networks~\cite{isola2017image,zhu2017toward,wang2018high,park2019semantic,choi2018stargan}. Additionally, diffusion models have been successfully applied to image-to-image translation tasks~\cite{saharia2022palette,voynov2022sketch,wang2022pretraining}. Recent advancements include ControlNet~\cite{zhang2023adding}, which enables control over pretrained diffusion models by finetuning and integrating supplementary control signals, thus achieving state-of-the-art controllable generation outcomes. In parallel, T2I-Adapter~\cite{mou2023t2i} also improves controllability by finetuning pretrained diffusion models. In line with the experimental settings of \cite{zhang2023adding,mou2023t2i}, we employ OFT to finetune pretrained diffusion models and consistently attain enhanced controllability, even when using less training data and fewer finetuned parameters. Importantly, OFT does not introduce additional computation during inference at test time.

\textbf{Model finetuning}. Adapting large pretrained models to specific downstream tasks through finetuning has gained substantial traction recently~\cite{devlin2018bert,brown2020language,he2022masked}. Within this context, various adaptation techniques (\emph{e.g.}, \cite{hulora2022,houlsby2019parameter,pfeiffer2020adapterfusion}) have been extensively explored, especially in the field of natural language processing. OFT shares similarities with LoRA~\cite{hulora2022}, which incorporates a low-rank assumption for weight updates during finetuning. However, OFT distinguishes itself by employing a layer-wise orthogonal transformation that updates neuron weights multiplicatively rather than additively. This approach rigorously maintains the pairwise angular relationships among neurons within a layer, contributing to improved training stability.

\section{Orthogonal Finetuning}

\subsection{Why Does Orthogonal Transformation Make Sense?}

To clarify the motivation for employing orthogonal transformation when finetuning text-to-image diffusion models, we analyze this issue through two subquestions: (1) the rationale for adjusting neuron angles (\emph{i.e.}, their directional attributes), and (2) the justification for using orthogonal transformations specifically for tuning these angles.

Addressing the first question, we take cues from empirical findings reported in \cite{liu2018decoupled,chen2020angular}, which suggest that angular differences in features effectively capture the semantic disparity. SphereNet~\cite{liu2017deep} demonstrates that enforcing normalization of all neuron outputs onto a unit hypersphere during neural network training results in comparable model capacity and potentially enhanced generalization, indicating that neuron directionality preserves the most significant data attributes. To further illustrate the role of neuron orientation, we devise a simplified experiment, depicted in Figure~\ref{fig:toy_ae}, training a conventional convolutional autoencoder on flower image data. Throughout training, the standard inner product generates the feature map, with $z$ as the convolution kernel output ($\bm{w}$ as the kernel and $\bm{x}$ as the sliding window input). In the test phase, we explore three alternative approaches for obtaining the feature map: (a) reusing the inner product, (b) extracting magnitude components, and (c) utilizing only angular components. The findings in Figure~\ref{fig:toy_ae} highlight that neuron angular data alone can nearly reconstruct the input images entirely, whereas magnitude values are largely uninformative. It is important to note that the cosine activation (c) is excluded during training; training relies solely on the inner product. This outcome underscores the predominance of neuron direction in encapsulating the semantic content of input images. Hence, adjusting neuron angles is likely a more effective strategy for modifying image semantics.

For the second question, a straightforward method for adjusting neuron orientation is to rotate or reflect all neurons within a given layer together, which aligns with employing orthogonal transformations. While one could consider alternative angular manipulations—such as individually rotating each neuron by varying amounts—our experiments indicate that orthogonal transformations strike an optimal balance between expressiveness and structural regularity. Additionally, \cite{liu2021orthogonal} provides evidence for the adequacy of orthogonal transformations in training neural networks. To further substantiate our position, we conduct an experiment focusing on the regularizing influence of orthogonality. Using the ControlNet~\cite{zhang2023adding} experimental setup for controllable generation, as visualized in Figure~\ref{fig:ortho}, we observe that the default OFT method achieves stable performance and precise control by the end of training (epoch 20). Conversely, removing the orthogonality constraint from OFT prevents the network from producing realistic images or exhibiting control. These observations confirm the essential role played by the orthogonality constraint in OFT.

\subsection{General Framework}

In the standard approach to finetuning, gradient descent is applied to the entire model or selected layers, generally utilizing a small learning rate. This restrained step size serves to implicitly restrict the deviation from the pretrained parameters, such that finetuning effectively minimizes $\thickmuskip=2mu \medmuskip=2mu \| \bm{M} - \bm{M}^0\|$, where $\bm{M}$ and $\bm{M}^0$ represent the finetuned and original pretrained weights, respectively. While this inherent regularization aligns with intuition, it may nevertheless grant excessive flexibility when adjusting large models. To address this, LoRA introduces a further low-rank restriction on the weight modification: $\thickmuskip=2mu \medmuskip=2mu \text{rank}(\bm{M} - \bm{M}^0)=r'$, with $r'$ being a predetermined small integer. In contrast, OFT imposes a different form of constraint, preserving pairwise neuron relationships by enforcing $\thickmuskip=2mu \medmuskip=2mu \| \text{HE}(\bm{M})-\text{HE}(\bm{M}^0) \|=0$, where $\text{HE}(\cdot)$ stands for the hyperspherical energy associated with the weight matrix.

For clarification, consider a fully connected layer expressed as $\thickmuskip=2mu \medmuskip=2mu \bm{W}=\{\bm{w}_1,\cdots,\bm{w}_n\}\in\mathbb{R}^{d\times n}$, with the $i$-th column $\bm{w}_i\in\mathbb{R}^{d}$ representing an individual neuron (with pretrained counterpart $\bm{W}^0$). The output $\thickmuskip=2mu \medmuskip=2mu \bm{z}\in\mathbb{R}^n$ for an input vector $\thickmuskip=2mu \medmuskip=2mu \bm{x}\in\mathbb{R}^d$ is computed as $\thickmuskip=2mu \medmuskip=2mu \bm{z}=\bm{W}^\top\bm{x}$. OFT interprets the finetuning process as one that seeks to minimize the gap in hyperspherical energy between the adjusted and original weights:

\begin{equation}\label{eq:oft_principle}
\footnotesize
\min_{\bm{W}} \left\| \text{HE}(\bm{W})-\text{HE}(\bm{W}^0)\right\|~~\Leftrightarrow~~\min_{\bm{W}} \bigg{\|} \sum_{i\neq j} \|\hat{\bm{w}}_i-\hat{\bm{w}}_j\|^{-1}-\sum_{i\neq j} \|\hat{\bm{w}}^0_i-\hat{\bm{w}}^0_j\|^{-1}\bigg{\|}
\end{equation}

Here, $\thickmuskip=2mu \medmuskip=2mu \hat{\bm{w}}_i:=\bm{w}_i/\|\bm{w}_i\|$ normalizes the $i$-th neuron, and the hyperspherical energy for $\bm{W}$ is given by $\thickmuskip=2mu \medmuskip=2mu \text{HE}(\bm{W}):= \sum_{i\neq j} \|\hat{\bm{w}}_i-\hat{\bm{w}}_j\|^{-1}$. Notably, Eq.~\eqref{eq:oft_principle} reaches a minimal value of zero, which is attainable provided that $\bm{W}$ can be obtained from $\bm{W}^0$ by an orthogonal transformation; that is, $\thickmuskip=2mu \medmuskip=2mu \bm{W}=\bm{R}\bm{W}^0$ where $\thickmuskip=2mu \medmuskip=2mu \bm{R}\in\mathbb{R}^{d\times d}$ is an orthogonal matrix (with determinant $1$ or $-1$ indicating a rotation or a reflection, respectively). The principal insight behind OFT is that model adaptation can be effectively achieved by learning a shared orthogonal matrix per layer, which transforms all neurons collectively. More concretely, for a given pretrained fully connected layer

\begin{equation}\label{eq:oft_general}
    \footnotesize
    \bm{z}=\bm{W}^\top\bm{x}=(\bm{R}\cdot\bm{W}^0)^\top\bm{x},~~~\text{s.t.}~\bm{R}^\top\bm{R}=\bm{R}\bm{R}^\top=\bm{I}
\end{equation}

In this formulation, $\bm{W}$ stands for the weight matrix obtained after OFT-finetuning, and $\bm{I}$ indicates the identity matrix. The framework of OFT is depicted in Figure~\ref{fig:block}. Mirroring the strategy of zero initialization employed in LoRA, it is essential that OFT allows finetuning from the original pretrained parameters $\bm{W}^0$. This requirement is fulfilled by initializing the orthogonal matrix $\bm{R}$ as an identity matrix, thereby ensuring that the finetuning process begins with the pretrained values. To preserve the orthogonality constraint on $\bm{R}$, one may leverage differentiable orthogonalization methods described in \cite{liu2021orthogonal,lezcano2019cheap}. Subsequently, we elaborate on efficient mechanisms for maintaining matrix orthogonality.

\subsection{Efficient Orthogonal Parameterization}\label{sect:eff_ortho}

While classical approaches like the Gram-Schmidt process provide differentiable orthogonalization, they are computationally intensive and may not be practical for large-scale scenarios~\cite{liu2021orthogonal}. To circumvent this inefficiency, we utilize the Cayley parameterization for generating orthogonal matrices. Concretely, the orthogonal matrix is defined as $\thickmuskip=2mu \medmuskip=2mu \bm{R}=(\bm{I}+\bm{Q})(I-\bm{Q})^{-1}$, where $\bm{Q}$ is a skew-symmetric matrix such that $\thickmuskip=2mu \medmuskip=2mu \bm{Q}=-\bm{Q}^\top$. This approach, while computationally favorable, confines the resultant orthogonal matrices to have determinant $1$, restricting $\bm{R}$ to the special orthogonal group. Empirically, this restriction does not impede model performance. Despite the use of Cayley transformation, parameterizing $\bm{R}$ for large $d$ remains parameter-intensive. To mitigate this, we propose adopting a block-diagonal structure for $\bm{R}$ with $r$ constituent blocks, yielding the representation:

\begin{equation}
    \footnotesize
    \bm{R}=\text{diag}(\bm{R}_1,\bm{R}_2,\cdots,\bm{R}_r)=\begin{bmatrix}
    \bm{R}_1\in \text{O}(\frac{d}{r}) &  &\\
    &\ddots & \\
    & & \bm{R}_r\in \text{O}(\frac{d}{r})
    \end{bmatrix}\in \text{O}(d)
\end{equation}

Here, $\text{O}(d)$ refers to the orthogonal group of degree $d$. We denote $\thickmuskip=2mu \medmuskip=2mu \bm{R}\in\mathbb{R}^{d\times d}$ as a general orthogonal matrix, while $\thickmuskip=2mu \medmuskip=2mu \bm{R}_i\in\mathbb{R}^{d/r\times d/r}$ (for all $i$) represent the block orthogonal matrices. When we set $\thickmuskip=2mu \medmuskip=2mu r=1$, the block-diagonal structure reduces to a standard full orthogonal matrix without block restrictions. For an orthogonal matrix of dimension $\thickmuskip=2mu \medmuskip=2mu d\times d$, the number of free parameters equals $\thickmuskip=2mu \medmuskip=2mu d(d-1)/2$, which yields a complexity of $\mathcal{O}(d^2)$. Conversely, when employing an $r$-block diagonal orthogonal matrix, the parameter count is $\thickmuskip=2mu \medmuskip=2mu d(d/r-1)/2$, resulting in $\thickmuskip=2mu \medmuskip=2mu \mathcal{O}(d^2/r)$ complexity. Furthermore, one may tie parameters across different blocks—specifically, set $\thickmuskip=2mu \medmuskip=2mu \bm{R}_i = \bm{R}_j$ for any $i\neq j$—which further brings the number of parameters down to $\mathcal{O}(d^2/r^2)$. These parameter-reduction techniques do not compromise the orthogonality of $\bm{R}$; thus, the preservation of hyperspherical energy is maintained.

We next contrast the parameter efficiency of OFT and LoRA. For LoRA with a low-rank setting $r'$, the total number of parameters to train becomes $\thickmuskip=2mu \medmuskip=2mu r'(d+n)$. If both $r$ and $r'$ are chosen proportional to $d$, for example $\thickmuskip=2mu \medmuskip=2mu r=r'=\alpha d$ where $\thickmuskip=2mu \medmuskip=2mu 0<\alpha\leq 1$ is a constant, then LoRA requires $\mathcal{O}(d^2+dn)$ parameters, while the OFT setup yields a much lower complexity of $\mathcal{O}(d)$. To make this comparison concrete, consider a weight matrix sized $\thickmuskip=2mu \medmuskip=2mu 128\times 128$; with $\thickmuskip=2mu \medmuskip=2mu r'=8$, LoRA results in $2,048$ parameters to train, whereas OFT needs only $960$ parameters if $\thickmuskip=2mu \medmuskip=2mu r=8$, and no block sharing is employed.

\subsection{Constrained Orthogonal Finetuning}

The expressive capacity of OFT can be further restricted by requiring the finetuned weights to remain close to their pretrained counterparts. In this setting, COFT adopts the following functional form:

\begin{equation}\label{eq:coft_general}
    \footnotesize
    \bm{z}=\bm{W}^\top\bm{x}=(\bm{R}\cdot\bm{W}^0)^\top\bm{x},~~~\text{s.t.}~\bm{R}^\top\bm{R}=\bm{R}\bm{R}^\top=\bm{I},~\norm{\bm{R}-\bm{I}}\leq \epsilon
\end{equation}

Here, we not only impose that $\bm{R}$ be orthogonal, but also restrict its deviation from the identity matrix by at most $\epsilon$. The Cayley parameterization from Section~\ref{sect:eff_ortho} can handle the orthogonality requirement. However, incorporating the $\epsilon$-deviation limit into this Cayley form is nontrivial. To build further understanding of the Cayley transform, we utilize the Neumann series to approximate $\thickmuskip=2mu \medmuskip=2mu \bm{R}=(\bm{I}+\bm{Q})(I-\bm{Q})^{-1}$

series converges in the operator norm). Consequently, we are able to reformulate the constraint $\thickmuskip=2mu \medmuskip=2mu\|\bm{R}-\bm{I}\|\leq\epsilon$ within the Cayley transform itself. The resulting equivalent formulation becomes $\thickmuskip=2mu \medmuskip=2mu \|\bm{Q}-\bm{0}\|\leq\epsilon'$, where $\bm{0}$ indicates the zero matrix and $\epsilon'$ denotes an alternative error tolerance parameter, distinct from $\epsilon$. This updated restriction on $\bm{Q}$ can be straightforwardly handled through projected gradient descent. To initialize the orthogonal matrix $\bm{R}$ as the identity, we set $\bm{Q}$ initially to the zero matrix. Conceptually, COFT imposes two explicit conditions: a small hyperspherical energy deviation and a bounded divergence from the pretrained weights. Whereas many existing finetuning approaches treat the second criterion as implicit, COFT formalizes it explicitly. Although OFT performs admirably, our empirical observations indicate that COFT delivers enhanced stability during finetuning, attributable to this explicit deviation constraint. Figure~\ref{fig:eps_coft} illustrates how varying $\epsilon$ impacts COFT's outcomes. The figure demonstrates that $\epsilon$ regulates the degree of adaptability during finetuning: increasing $\epsilon$ leads COFT's output to converge towards that of OFT, while reducing $\epsilon$ constrains the COFT results to closely mirror those of the original pretrained text-to-image diffusion model.

\subsection{Re-scaled Orthogonal Finetuning}

We introduce a straightforward augmentation to the standard OFT framework by incorporating a learnable scaling factor per neuron. This modification arises from the observation that scaling does not affect hyperspherical energy, since normalization removes any change in magnitude. Our approach employs the following forward computation: $\thickmuskip=2mu \medmuskip=2mu\bm{z}=(\bm{R}\bm{W}^0\bm{D})^\top\bm{x}$\footnote{\textbf{Errata}: In the NeurIPS camera-ready submission, the forward pass for re-scaled OFT was incorrectly stated as $\thickmuskip=2mu \medmuskip=2mu\bm{z}=(\bm{D}\bm{R}\bm{W}^0)^\top\bm{x}$. The code implementation was correct, so findings in Appendix~\ref{app:rescaled} are valid.}, where $\thickmuskip=2mu \medmuskip=2mu \bm{D}=\text{diag}(s_1,\cdots,s_n)\in\mathbb{R}^{n\times n}$ is a diagonal matrix with learnable positive entries $s_1,\ldots,s_n$. While the original OFT forward pass in Eq.~\eqref{eq:oft_general} optimizes only $\bm{R}$, our formulation treats both $\bm{D}$ and $\bm{R}$ as trainable parameters. This re-scaled OFT variant enhances OFT’s expressive power at the cost of only a trivial increase in parameter count. To highlight the standalone effect of orthogonal transformations, our experiments mainly use the original OFT, although in general the re-scaled variant yields superior performance (see Appendix~\ref{app:rescaled}).

\section{Intriguing Insights and Discussions}

\textbf{OFT’s generality across architectures}. In principle, OFT can be adopted for any neural network architecture. For example, while LoRA is typically restricted to the attention weights of Transformers~\cite{hulora2022}, our experiments limit OFT to attention weights as well for a fair comparison. Notably, OFT is particularly suitable for convolutional layers due to the meaningful interpretation of the block-diagonal structure of $\bm{R}$, which LoRA does not naturally provide. By aligning the number of blocks in $\bm{R}$ with the number of input channels, each block can independently transform a single neuron channel—analogous to learning depth-wise convolutional filters~\cite{chollet2017xception}. Conversely, if all blocks in $\bm{R}$ are identical, the orthogonal transformation is shared across all channels. We present initial experiments with OFT on convolutional layers in Appendix~\ref{app:convolution}.

\textbf{Connection to LoRA}. LoRA mitigates significant deviations in the pretrained weight matrix by introducing a low-rank update, thus preserving its original information. In comparison, OFT imposes an orthogonal (full-rank) constraint on the transformation applied to the pretrained weights, ensuring the pretraining knowledge remains intact. The forward computation in OFT can be reformulated as $\thickmuskip=2mu \medmuskip=2mu \bm{z}=(\bm{R}\bm{W}^0)^\top\bm{x}=(\bm{W}^0+(\bm{R}-\bm{I})\bm{W}^0)^\top\bm{x}$, where $\thickmuskip=2mu \medmuskip=2mu (\bm{R}-\bm{I})\bm{W}^0$ acts similarly to LoRA's low-rank update to the weights. Because $\bm{W}^0$ is generally full-rank, OFT also induces a low-rank update when $\thickmuskip=2mu \medmuskip=2mu \bm{R}-\bm{I}$ is constrained to be low-rank. Just as LoRA utilizes a rank parameter $r'$ to decrease trainable parameters, OFT incorporates a diagonal block size parameter $r$ to achieve parameter reduction. Notably, while both methods enhance parameter efficiency, they rely on different underlying strategies: LoRA exploits the low-rank nature of updates, whereas OFT capitalizes on sparse structure via block-diagonal orthogonal constraints.

\textbf{Why OFT converges faster}. Evidence from Figure~\ref{fig:toy_ae} indicates that changing the orientation of neurons offers the most substantial impact on the encoded semantics, which is precisely the focus of OFT. Furthermore, the OFT method can be interpreted as conducting finetuning on the hypersphere manifold of neurons, resulting in a smoother optimization landscape and improved training dynamics. This theoretical advantage is supported by empirical results presented in \cite{liu2021orthogonal}.

\textbf{Why not minimize hyperspherical energy}. Our approach diverges from that of \cite{liu2021orthogonal} by not pursuing hyperspherical energy minimization. In classification settings, it is advantageous for neurons to avoid redundancy. Achieving minimum hyperspherical energy causes neurons to be distributed uniformly across the hypersphere, which, while optimal for de novo feature design, is not suitable for finetuning as it may erase beneficial features acquired during pretraining.

\textbf{Balancing flexibility and regularity in finetuning.} Our analysis reveals a fundamental compromise between model adaptability and imposed structure. While conventional finetuning achieves maximal adaptability, it often comes at the expense of stability and frequently leads to model collapse. In contrast, OFT employs a notably straightforward approach that maintains an effective trade-off between flexibility and regularity, achieved by conserving the pairwise angles among neurons. Additionally, employing block-diagonal parameterization serves to further constrain the orthogonal transformation, introducing a heightened degree of regularization.

\textbf{No inference-time computational cost.} In contrast to ControlNet, the OFT approach does not impose any computational burden during inference. Specifically, during inference, one can efficiently compute the resulting weight matrix by directly left-multiplying the learned orthogonal matrix $\bm{R}$ to the original pretrained weights $\bm{W}^0$, resulting in the equivalent parameterization $\thickmuskip=2mu \medmuskip=2mu \bm{W}=\bm{R}\bm{W}^0$. This approach ensures that inference proceeds at the same speed as when using the original pretrained model.

\section{Experiments and Results}

\textbf{General protocol.} For all experiments, we utilize Stable Diffusion v1.5~\cite{rombach2022high} as our backbone text-to-image model. To ensure an unbiased comparison, image samples generated by each method are selected randomly. Subject-driven generation tasks primarily follow the protocol of DreamBooth~\cite{ruiz2023dreambooth}, whereas experiments concerning controllable generation reference the practices in ControlNet~\cite{zhang2023adding} and T2I-Adapter~\cite{mou2023t2i}. For fairness in benchmarking against LoRA, both OFT and COFT are applied exclusively to the layers where LoRA modifications occur. A comprehensive presentation of additional results and detailed experimental settings can be found in Appendix~\ref{app:settings}.

\subsection{Subject-driven Generation}

\textbf{Experimental configuration.} DreamBooth~\cite{ruiz2023dreambooth} and LoRA~\cite{hulora2022} serve as comparative baselines in this study. All techniques leverage the same loss formulation adopted by DreamBooth. For both DreamBooth and LoRA, we adhere closely to the methodologies and optimal hyperparameter settings recommended in their respective original works. Additional findings are documented in Appendix~\ref{app:settings},\ref{app:coft_oft},\ref{app:more_qualitative},\ref{app:failure}.

\textbf{Finetuning stability and convergence}. Our initial analysis focuses on the stability and rate of convergence during finetuning for DreamBooth, LoRA, OFT, and COFT. The relevant results are depicted in Figure~\ref{fig:teaser} and Figure~\ref{fig:db_convergence}. From these figures, it is evident that both COFT and OFT demonstrate robust and stable behavior when finetuning the diffusion model. At 400 training steps, DreamBooth and both OFT variants provide satisfactory control over the generated outputs, whereas LoRA struggles to maintain the integrity of subject identity. Advancing to 2000 steps, DreamBooth begins to exhibit mode collapse, generating degenerate outputs, and LoRA fails to accurately render the yellow shirt (opting instead for yellow fur). In sharp contrast, OFT and COFT consistently produce stable and well-controlled outputs throughout training. These findings corroborate the rapid and stable convergence properties of the OFT approach for subject-driven synthesis. The insensitivity to the exact number of finetuning steps emerges as a significant practical advantage, reducing the burden of tailoring the iteration count for individual subjects. With both OFT and COFT, one can select a sufficiently high iteration count without engaging in meticulous parameter tuning. When COFT is combined with an appropriate choice of $\epsilon$, tuning of both the learning rate and the number of iterations becomes trivial.

\textbf{Quantitative comparison}. In alignment with the protocol of \cite{ruiz2023dreambooth}, we quantitatively assess subject fidelity (via DINO~\cite{caron2021emerging} and CLIP-I~\cite{radford2021learning}), prompt fidelity (using CLIP-T~\cite{radford2021learning}), and diversity of generated samples (measured by LPIPS~\cite{zhang2018unreasonable}). For CLIP-I, we calculate the mean pairwise cosine similarity between the CLIP embeddings of synthesized and authentic images; DINO is analogous, substituting ViT S/16 DINO features for the CLIP representation. The CLIP-T metric is computed as the mean cosine similarity between text prompt embeddings and the respective image embeddings. We further assess diversity by evaluating the average LPIPS cosine similarity between pairs of generated images depicting the same subject and conditioned on the same prompt. As reported in Table~\ref{table:dreambooth_final}, COFT and OFT exceed the performance of DreamBooth and LoRA by a notable margin on DINO and CLIP-I, while maintaining marginally better or similar results on the prompt fidelity and diversity benchmarks. To ensure reliability, we repeated finetuning for each method 30 times using different random seeds with the same pretrained model, thereby reducing the effect of stochasticity. These outcomes highlight that OFT not only confers improvements in convergence and robustness but also consistently delivers superior final outcomes.

\textbf{Qualitative comparison}. To provide a clearer illustration of the advantages offered by OFT, Figure~\ref{fig:dreambooth_final} presents a selection of randomly chosen samples for subject-driven image generation. To ensure an unbiased evaluation, all samples originate from the same finetuned model under each respective method, and no additional text prompt optimization is performed for any specific result. For each approach, we utilize the checkpoint that attains the highest validation CLIP scores. As illustrated in Figure~\ref{fig:dreambooth_final}, OFT and COFT both demonstrate strong capabilities in retaining semantic subject characteristics, in contrast to LoRA, which frequently does not maintain the integrity of the subject's identity (\emph{e.g.}, in the bowl sample, LoRA completely fails to preserve the subject). Additionally, OFT and COFT display substantially superior adherence to text prompt guidance, while DreamBooth, though effective at maintaining subject identity, often struggles to align the generated images with the provided textual prompts (\emph{e.g.}, see the initial row for the bowl example). Overall, this comparison underscores that OFT excels in balancing both controllability and the retention of subject information. Furthermore, OFT exhibits robustness with respect to the number of iterations, delivering consistent performance even with many iterations—something that neither DreamBooth nor LoRA manages effectively. Supplementary qualitative examples can be found in Appendix~\ref{app:more_qualitative}. Additionally, the human evaluation presented in Appendix~\ref{app:human} offers further evidence in support of OFT’s advantages.

\subsection{Controllable Generation}

\textbf{Settings}. For baseline comparison, we employ ControlNet~\cite{zhang2023adding}, T2I-Adapter~\cite{mou2023t2i}, and LoRA~\cite{hulora2022}. The main text considers three demanding controllable generation scenarios: Canny edge to image~(C2I) using the COCO dataset~\cite{lin2014microsoft}, segmentation map to image~(S2I) on the ADE20K dataset~\cite{zhou2017scene}, and landmark to face~(L2F) with the CelebA-HQ dataset~\cite{karras2018progressive,xia2021tedigan}. All methods perform finetuning on Stable Diffusion~(SD) v1.5 for 20 epochs across these three datasets. Additional results are available in Appendix~\ref{app:more_qualitative},~\ref{app:more_control_task}, and~\ref{app:failure}.

\textbf{Convergence}. We assess how quickly ControlNet, T2I-Adapter, LoRA, and COFT converge when applied to the L2F task, presenting both numerical and visual analyses. As a measure of performance, we calculate the average $\ell_2$ distance between the landmarks from the control faces and those predicted by the models. Figure~\ref{fig:face_convergence} illustrates how the landmark error changes as models are fine-tuned over various epoch counts. Both COFT and OFT demonstrate markedly accelerated convergence rates. Notably, LoRA requires 20 epochs to attain the performance that OFT reaches at just 8 epochs. OFT and COFT involve a comparable number of trainable parameters to LoRA—which is substantially lower than that of ControlNet—yet they converge with greater efficiency compared to other existing approaches. Additionally, the efficiency of OFT's convergence is further substantiated by results in Figure~\ref{fig:teaser}, where the right panel highlights OFT’s superior data efficiency: it performs well with as little as 5\% of the ADE20K dataset, outpacing both ControlNet and LoRA. For visual comparison, we emphasize OFT, COFT, and ControlNet, given that ControlNet achieves landmark errors closest to those of OFT. According to Figure~\ref{fig:face_convergence_qual}, OFT and COFT display stable convergence, with the synthesized face poses incrementally matching the specified landmarks. In contrast, ControlNet exhibits a sudden leap in controllability around the 8th epoch, mirroring the abrupt convergence described in \cite{zhang2023adding}. The more gradual and reliable convergence of OFT facilitates practical application, as its training process is notably smoother and more foreseeable, thereby simplifying the tuning of hyperparameters for OFT.

\textbf{Quantitative comparison}. To assess the effectiveness of controllable generation, we propose a control consistency metric. This approach involves extracting the control signal from the generated output and measuring its alignment with the original input control. Specifically, for the C2I task, we report IoU and F1 score; for the S2I task, we measure mean IoU as well as mean and overall accuracy; for the L2F task, we calculate the average $\ell_2$ distance between the control and predicted landmarks. Additional information on these consistency measures is presented in Appendix~\ref{app:settings}. Across all compared approaches, hyperparameters are optimized to the best of our ability. According to the results in Table~\ref{table:control_gen}, OFT and COFT deliver significantly improved and more precise control relative to baseline methods. We note that adapter-based techniques (\emph{e.g.}, T2I-Adapter and ControlNet) demonstrate slower convergence and poorer final performance. In the S2I task, LoRA surpasses ControlNet, but it underperforms in C2I and L2F tasks. Generally, we observe a relatively limited upper-bound in LoRA's results, despite extensive hyperparameter tuning. In contrast, the OFT framework’s performance continues to improve as more parameters are made trainable, suggesting its potential is not yet fully realized. Our quantitative evaluation framework for controllable generation, capable of precisely benchmarking the control ability of finetuned models (to the extent permitted by pre-existing segmentation/detection tools), constitutes a novel contribution of this work.

\textbf{Qualitative comparison}. We further present a qualitative evaluation of OFT and COFT alongside prominent existing approaches such as ControlNet, T2I-Adapter, and LoRA. As depicted in Figure~\ref{fig:control_final}, images generated by OFT and COFT exhibit not only superior realism and high fidelity but also precise controllability. For the S2I task, LoRA fails to adhere to the structure provided by the input segmentation map, whereas ControlNet, OFT, and COFT are successful in imposing effective control over the outputs. OFT and COFT surpass ControlNet by generating images that feature richer details and more plausible geometric coherence, all while utilizing significantly fewer model parameters. Regarding the C2I task, OFT and COFT are capable of producing semantically consistent images given only rudimentary Canny edge maps, whereas T2I-Adapter and LoRA demonstrate considerably inferior performance. When it comes to the L2F task, our framework is notably superior at achieving precise pose guidance for facial images, even under difficult orientations. Across all three control scenarios, both OFT and COFT consistently deliver higher-quality results compared to current state-of-the-art baselines, underscoring the robustness of the OFT framework for controllable image generation. Additional qualitative examples covering these control tasks can be found in Appendix~\ref{app:control_exp}, and we further illustrate OFT’s capacity for a wider array of control challenges (such as dense pose to human body, sketch to image, and depth to image conversions) in Appendix~\ref{app:more_control_task}.

\section{Concluding Remarks and Open Problems}

Inspired by the critical role that angular relationships among neurons play in defining visual semantics, we introduce a straightforward yet powerful finetuning strategy—orthogonal finetuning—for controlling text-to-image diffusion models. Our approach is specifically designed to address two primary tasks: subject-driven generation and controllable generation. In comparison to prior approaches, OFT achieves superior control and greater stability during finetuning, while also requiring a reduced number of tunable parameters. Furthermore, OFT avoids additional inference cost, thereby facilitating practical deployment.

The adoption of OFT also brings forth several noteworthy unresolved questions. Firstly, OFT enforces orthogonality using Cayley parametrization, which necessitates computing a matrix inverse and can marginally hinder scalability. While our use of block diagonal parametrization alleviates this constraint to some extent, devising methods to efficiently compute the matrix inverse in a differentiable manner remains an open issue. Secondly, the compositional potential of OFT is notable: orthogonal matrices derived from separate OFT finetuning processes can be combined through multiplication and will retain orthogonality. Investigating whether this composition preserves information from each individual downstream task is a promising direction for future work. Lastly, the parameter efficiency of OFT heavily relies on the use of block diagonal structures, which may inherently introduce certain biases and restrict flexibility. Identifying methods to further enhance parameter efficiency with reduced bias and greater effectiveness constitutes a significant and ongoing research challenge.

\end{document}